
Sedmnáctý případ užití je specifikován v tabulce č. \ref{tab:usecase:UC17}.

\begin{UseCaseTable}
  {Stažení archivu naměřených dat}
  {UC17}
  {F17}
  {Získání periodicky aktualizovaného archivu všech dosud naměřených dat nebo jeho vymezené podmnožiny pro hromadné zpracování}
  {Výzkumný pracovník}
  {}
  {Žádné -- archiv je veřejně dostupný}
  {Konzument získal archiv naměřených dat v otevřeném formátu}

\UseCaseScenario{Základní scénář -- stažení úplného archivu}
\UseCaseStep{výzkumný pracovník}{Odešle systému HTTP požadavek na endpoint poskytující archiv všech dosud naměřených dat.}
\UseCaseStep{systém}{Vrátí periodicky aktualizovaný archiv v otevřeném strojově čitelném formátu.}

\UseCaseScenario{Alternativní scénář -- stažení vymezené podmnožiny archivu}
\UseCaseStepCustom{A1}{výzkumný pracovník}{Odešle systému požadavek na archiv s parametry vymezujícími podmnožinu (časový interval, geografická oblast, identifikátory konkrétních entit).}
\UseCaseStepCustom{A2}{systém}{Ověří, že zadané parametry jsou v očekávaném tvaru.}
\UseCaseStepCustom{A3}{systém}{Vrátí archiv obsahující pouze data odpovídající zadané podmnožině.}

\UseCaseScenario{Alternativní scénář -- neplatné parametry vymezení}
\UseCaseStepCustom{B1}{výzkumný pracovník}{shodné s krokem A1 alternativního scénáře.}
\UseCaseStepCustom{B2}{systém}{Zjistí, že některý z parametrů vymezení má neplatný formát nebo odkazuje na neexistující entitu.}
\UseCaseStepCustom{B3}{systém}{Vrátí chybovou odpověď s odpovídajícím HTTP stavovým kódem a popisem problematického parametru.}

\end{UseCaseTable}
